\section{Базовые конструкции C++}

\subsection{Алфавит языка C++. Простая программа C++, структура, основные действующие лица}
\begin{definition}[Алфавит языка]
Алфавит C++ — это набор допустимых символов, из которых строится исходный текст программы: латинские буквы, цифры, знаки препинания, пробелы, табуляции, символы переноса строки и специальные символы препроцессора.
\end{definition}

В исходном тексте программы обычно выделяют следующие «действующие лица»:
\begin{itemize}
    \item \textbf{препроцессор} — обрабатывает директивы \texttt{\#include}, \texttt{\#define}, \texttt{\#if} и т.д.;
    \item \textbf{компилятор} — проверяет синтаксис и семантику, строит объектный код;
    \item \textbf{линкер} — собирает объектные файлы и библиотеки в исполняемую программу;
    \item \textbf{runtime-среда} — окружение, в котором программа реально выполняется.
\end{itemize}

Минимальная программа в C++ обычно содержит функцию \texttt{main}:
\begin{verbatim}
int main() {
    return 0;
}
\end{verbatim}

\begin{remark}
С точки зрения языка C++ это уже полноценная программа. Всё остальное — подключение библиотек, объявление типов, функции, классы — лишь расширяет её структуру.
\end{remark}

\subsection{Ключевые слова. Основные ключевые слова}
\begin{definition}[Ключевое слово]
\textit{Keyword} — это зарезервированное слово языка C++, имеющее специальный синтаксический смысл и не допускающееся в качестве идентификатора.
\end{definition}

К основным ключевым словам относятся:
\begin{itemize}
    \item для типов и объявлений: \texttt{int}, \texttt{double}, \texttt{char}, \texttt{bool}, \texttt{class}, \texttt{struct}, \texttt{enum}, \texttt{void};
    \item для управления потоком: \texttt{if}, \texttt{else}, \texttt{switch}, \texttt{case}, \texttt{for}, \texttt{while}, \texttt{do}, \texttt{break}, \texttt{continue}, \texttt{return};
    \item для описания сущностей: \texttt{const}, \texttt{constexpr}, \texttt{static}, \texttt{extern}, \texttt{inline}, \texttt{virtual}, \texttt{namespace};
    \item для препроцессора и специальных механизмов: \texttt{new}, \texttt{delete}, \texttt{try}, \texttt{catch}, \texttt{throw}, \texttt{public}, \texttt{private}, \texttt{protected}.
\end{itemize}

\begin{remark}
Не все «служебные слова» языка являются ключевыми в строгом смысле, но для коллоквиума важно понимать, что они имеют фиксированное значение и не могут использоваться произвольно.
\end{remark}

\subsection{Идентификаторы, правило составления идентификаторов. Влияние кодстайла на правило именования}
\begin{definition}[Идентификатор]
\textit{Identifier} — имя объекта, функции, типа, namespace, макроса и других сущностей программы.
\end{definition}

Формально идентификатор в C++:
\begin{itemize}
    \item состоит из букв, цифр и символа подчёркивания;
    \item не может начинаться с цифры;
    \item чувствителен к регистру;
    \item не должен совпадать с ключевыми словами.
\end{itemize}

\begin{remark}
Стандарт языка даёт лишь минимальные ограничения. Практика code style добавляет дополнительные соглашения: например, типы пишут в \texttt{PascalCase}, функции и переменные — в \texttt{camelCase}, константы — в \texttt{kPascalCase} или \texttt{UPPER\_CASE}, а макросы — только \texttt{UPPER\_CASE}.
\end{remark}

\subsection{Инструкция (statement) как предложение языка. Классификация инструкций}
\begin{definition}[Инструкция]
\textit{Statement} — это законченная синтаксическая конструкция языка C++, выражающая действие или описание.
\end{definition}

Обычно инструкции разделяют на несколько больших классов:
\begin{itemize}
    \item \textbf{инструкция-описание} — объявляет сущность;
    \item \textbf{инструкция-определение} — создаёт сущность;
    \item \textbf{инструкция-выражение} — вычисляет выражение и завершается \texttt{;};
    \item \textbf{инструкции управления потоком} — \texttt{if}, \texttt{switch}, \texttt{for}, \texttt{while}, \texttt{do}, \texttt{break}, \texttt{continue}, \texttt{return};
    \item \textbf{блок инструкций} — составная инструкция в фигурных скобках.
\end{itemize}

\begin{remark}
Не каждое законченное действие в тексте программы начинается с новой строки, и не каждая строка — это отдельная инструкция. Видимым признаком инструкции обычно служит точка с запятой или фигурные скобки, но есть исключения: например, блок и управляющие конструкции могут не оканчиваться \texttt{;}.
\end{remark}

\subsection{Инструкции описания объекта данных и объекта действия (функции). \texttt{extern}. Прототип функции}
\begin{definition}[Описание]
\textit{Declaration} — это сообщение компилятору о существовании сущности и её свойствах без обязательного создания самой сущности.
\end{definition}

\begin{definition}[Описание объекта данных]
Это инструкция, которая сообщает компилятору о переменной или константе:
\begin{verbatim}
extern int x;
const double pi = 3.14159;
\end{verbatim}
\end{definition}

\begin{definition}[Описание функции]
Это инструкция, которая сообщает сигнатуру функции:
\begin{verbatim}
int sum(int a, int b);
\end{verbatim}
Такая запись называется \textit{prototype} или \textit{function declaration}.
\end{definition}

Ключевое слово \texttt{extern} обычно означает, что определение находится в другом месте:
\begin{verbatim}
extern int counter;
\end{verbatim}

\begin{remark}
\texttt{extern} особенно важно для глобальных объектов и связки между несколькими translation units.
\end{remark}

\subsection{Инструкции определения и описания объекта данных}
\begin{definition}[Определение]
\textit{Definition} — это инструкция, которая создаёт сущность. Для объекта данных это обычно означает выделение памяти.
\end{definition}

\begin{example}
\begin{verbatim}
int a;          // definition + declaration
extern int b;   // declaration only
int c = 10;     // definition + declaration
\end{verbatim}
\end{example}

\begin{remark}
В C++ определение объекта данных почти всегда одновременно является и его описанием. Исключение — \texttt{extern}-объявление, которое только описывает объект.
\end{remark}

\subsection{Инициализация объектов. Конструкции \texttt{=}, \texttt{()}, \texttt{\{\}}, \texttt{=\{\}}}
\begin{definition}[Инициализация]
\textit{Initialization} — задание начального значения объекту в момент его создания.
\end{definition}

Основные формы:
\begin{itemize}
    \item \textbf{copy-initialization}: \texttt{T x = value;}
    \item \textbf{direct-initialization}: \texttt{T x(value);}
    \item \textbf{list-initialization}: \texttt{T x\{value\};}
    \item \textbf{copy-list-initialization}: \texttt{T x = \{value\};}
\end{itemize}

Общее:
\begin{itemize}
    \item все они могут задавать начальное значение;
    \item все относятся к созданию объекта.
\end{itemize}

Различия:
\begin{itemize}
    \item \texttt{\{\}} лучше ловит narrowing;
    \item \texttt{()} может вызывать конструкторы и иногда попадает в эффект \textit{most vexing parse};
    \item \texttt{=} исторически выглядит как присваивание, но при объявлении это именно инициализация;
    \item \texttt{=\{\}} сочетает синтаксис \texttt{=} и семантику списка инициализации.
\end{itemize}

\begin{example}
\begin{verbatim}
int a = 5;
int b(5);
int c{5};
int d = {5};
\end{verbatim}
\end{example}

\subsection{Блок инструкций. Scope}
\begin{definition}[Блок инструкций]
\textit{Compound statement} — это последовательность инструкций в фигурных скобках \texttt{\{} \texttt{\}}.
\end{definition}

Блок используется для:
\begin{itemize}
    \item последовательного выполнения нескольких инструкций;
    \item создания локальной области видимости;
    \item группировки инструкций в теле \texttt{if}, \texttt{for}, \texttt{while}, функции и т.д.
\end{itemize}

\begin{definition}[Scope]
\textit{Scope} — область видимости, в пределах которой имя доступно без дополнительной квалификации.
\end{definition}

Блок создаёт \textbf{local scope}: объявленные внутри него сущности живут и видны только до конца блока.

\begin{example}
\begin{verbatim}
{
    int x = 1;
    int y = x + 2;
}
\end{verbatim}
\end{example}

\subsection{Функция как именованный блок инструкций с заданным входом и выходом}
\begin{definition}[Функция]
\textit{Function} — именованный блок инструкций с параметрами на входе и, как правило, возвращаемым значением на выходе.
\end{definition}

Функция задаёт:
\begin{itemize}
    \item имя;
    \item список параметров;
    \item возвращаемый тип;
    \item тело — блок инструкций;
    \item контракт вызова.
\end{itemize}

\begin{example}
\begin{verbatim}
int sum(int a, int b) {
    return a + b;
}
\end{verbatim}
\end{example}

\subsection{Объекты и типы. Характеристики объектов и характеристик типов}
\begin{definition}[Объект]
Объект — это область памяти, в которой хранятся данные некоторого типа.
\end{definition}

\begin{definition}[Тип]
Тип — это описание множества значений и множества допустимых операций над ними.
\end{definition}

\textbf{Характеристики объекта}:
\begin{itemize}
    \item имя (может отсутствовать);
    \item тип;
    \item значение;
    \item адрес;
    \item размер;
    \item время жизни;
    \item область видимости имени.
\end{itemize}

\textbf{Характеристики типа}:
\begin{itemize}
    \item множество допустимых значений;
    \item операции;
    \item представление в памяти;
    \item размер и выравнивание;
    \item правила инициализации, преобразований и сравнения.
\end{itemize}

\begin{remark}
В C++ тип и объект тесно связаны, но это не одно и то же: тип определяет свойства, а объект является конкретной реализацией этих свойств в памяти.
\end{remark}

\subsection{Представление объектов в оперативной памяти. \texttt{sizeof()}}
\begin{definition}[Footprint]
\textit{Footprint} объекта — это то, сколько памяти объект занимает и как эта память устроена.
\end{definition}

На представление влияют:
\begin{itemize}
    \item тип объекта;
    \item размер типа;
    \item выравнивание (\textit{alignment});
    \item способ кодирования чисел;
    \item наличие служебных полей у классов;
    \item оптимизации компилятора.
\end{itemize}

Оператор \texttt{sizeof} позволяет узнать размер объекта или типа:
\begin{verbatim}
sizeof(int)
sizeof x
\end{verbatim}

\begin{itemize}
    \item применим к типам и объектам;
    \item результат имеет тип \texttt{std::size\_t};
    \item возвращает размер в байтах.
\end{itemize}

\begin{remark}
\texttt{sizeof} не вычисляет выражение как runtime-операция; обычно это compile-time информация.
\end{remark}

\subsection{Операционная семантика типа. Абстрактный тип данных}
\begin{definition}[Операционная семантика]
Операционная семантика типа — это набор операций, которые можно выполнять над объектами этого типа, и смысл этих операций.
\end{definition}

Например, для \texttt{int} доступны сложение, вычитание, сравнение, инкремент и т.д.

\begin{definition}[ADT]
\textit{Abstract Data Type} — абстрактный тип данных, для которого важны интерфейс и набор операций, а не конкретная внутренняя реализация.
\end{definition}

Примеры ADT:
\begin{itemize}
    \item стек;
    \item очередь;
    \item множество;
    \item ассоциативный словарь.
\end{itemize}

\subsection{Классификация типов в C++: трехтаксонная и двухтаксонная}
Классическая \textbf{трёхтаксонная} классификация:
\begin{itemize}
    \item \textbf{фундаментальные} типы;
    \item \textbf{производные} типы;
    \item \textbf{пользовательские} типы.
\end{itemize}

\begin{itemize}
    \item \textbf{Фундаментальные}: \texttt{int}, \texttt{char}, \texttt{bool}, \texttt{float}, \texttt{double}, \texttt{void}, \texttt{nullptr\_t}.
    \item \textbf{Производные}: указатели, ссылки, массивы, функции, указатели на члены.
    \item \textbf{Пользовательские}: \texttt{class}, \texttt{struct}, \texttt{enum}, \texttt{union}.
\end{itemize}

По cppreference часто используют \textbf{двухтаксонную} классификацию:
\begin{itemize}
    \item \textbf{fundamental};
    \item \textbf{compound (derived-like)}.
\end{itemize}

\begin{remark}
Для курса полезно уметь оперировать обеими схемами: трёхтаксонная удобнее для учебного объяснения, а двухтаксонная чаще встречается в документации.
\end{remark}

\subsection{Целочисленные типы данных: размеры, знаковость, диапазоны. \texttt{<cstdint>}, \texttt{<climits>}, \texttt{<limits>}, Little Endian}
\begin{definition}[Целочисленные типы]
К целочисленным типам относятся \texttt{char}, \texttt{short}, \texttt{int}, \texttt{long}, \texttt{long long} и их unsigned-версии.
\end{definition}

\begin{itemize}
    \item размеры зависят от платформы, кроме некоторых минимальных гарантий стандарта;
    \item знаковость определяется либо явно (\texttt{signed}/\texttt{unsigned}), либо по умолчанию;
    \item диапазон зависит от числа бит и способа кодирования.
\end{itemize}

Псевдонимы фиксированного размера из \texttt{<cstdint>}:
\begin{itemize}
    \item \texttt{std::int8\_t}, \texttt{std::uint8\_t};
    \item \texttt{std::int16\_t}, \texttt{std::uint16\_t};
    \item \texttt{std::int32\_t}, \texttt{std::uint32\_t};
    \item \texttt{std::int64\_t}, \texttt{std::uint64\_t}.
\end{itemize}

Диапазоны целых типов можно смотреть:
\begin{itemize}
    \item в \texttt{<climits>} — через макросы вроде \texttt{INT\_MIN}, \texttt{INT\_MAX};
    \item в \texttt{<limits>} — через \texttt{std::numeric\_limits<T>}.
\end{itemize}

\begin{remark}
Для хранения многобайтовых чисел важен порядок байт. На большинстве современных платформ используется \textit{Little Endian}: младший байт идёт по меньшему адресу.
\end{remark}

\subsection{Типы чисел с плавающей точкой: \texttt{float}, \texttt{double}}
Основные типы с плавающей точкой:
\begin{itemize}
    \item \texttt{float};
    \item \texttt{double};
    \item \texttt{long double}.
\end{itemize}

Их представление обычно основано на стандарте IEEE 754 и включает:
\begin{itemize}
    \item знак;
    \item порядок (exponent);
    \item мантиссу (fraction/significand).
\end{itemize}

\begin{itemize}
    \item \texttt{float} обычно занимает 4 байта;
    \item \texttt{double} обычно 8 байт;
    \item \texttt{long double} зависит от платформы.
\end{itemize}

Такие типы позволяют представлять:
\begin{itemize}
    \item обычные числа;
    \item очень большие и очень маленькие числа;
    \item специальные значения: \texttt{+inf}, \texttt{-inf}, \texttt{NaN}.
\end{itemize}

\subsection{Инициализация числовых объектов. Числовые литералы, narrowing, \texttt{=} и \texttt{\{\}}}
Числовой литерал может уточняться суффиксами:
\begin{itemize}
    \item \texttt{1u} — \texttt{unsigned};
    \item \texttt{1l} — \texttt{long};
    \item \texttt{1ll} — \texttt{long long};
    \item \texttt{1.0f} — \texttt{float};
    \item \texttt{1.0l} — \texttt{long double}.
\end{itemize}

\begin{definition}[Narrowing]
\textit{Narrowing} — сужение типа, при котором значение может потерять часть информации, например при переводе \texttt{double} в \texttt{int}.
\end{definition}

Сравнение:
\begin{itemize}
    \item \texttt{int x = 3.14;} — может пройти с потерей дробной части;
    \item \texttt{int x\{3.14\};} — ошибка компиляции из-за narrowing.
\end{itemize}

\begin{remark}
Именно поэтому список инициализации \texttt{\{\}} считается более безопасным.
\end{remark}

\subsection{POD-типы и их особенность}
\begin{definition}[POD]
\textit{POD} (\textit{Plain Old Data}) — исторический класс простых типов без сложного поведения. В современном C++ это понятие в значительной мере заменено более точными категориями вроде \textit{trivial} и \textit{standard-layout}.
\end{definition}

Примеры учебно:
\begin{itemize}
    \item \texttt{int};
    \item \texttt{double};
    \item простая \texttt{struct} без пользовательских конструкторов и виртуальных функций.
\end{itemize}

\begin{remark}
Особенность таких типов — простота представления в памяти и предсказуемое копирование побайтно, если это допускается стандартом.
\end{remark}

\subsection{Выражения и операторы. Арность, приоритет, instruction-expression. Инициализация и присваивание}
\begin{definition}[Выражение]
\textit{Expression} — конструкция языка, которая вычисляется и даёт значение или effect.
\end{definition}

\begin{definition}[Оператор и операнд]
Оператор — символ или ключевое слово, задающее операцию. Операнд — объект, над которым выполняется операция.
\end{definition}

\begin{itemize}
    \item по арности операторы бывают унарные, бинарные и тернарные;
    \item приоритет определяет порядок вычисления;
    \item если порядок неочевиден, его задают скобками.
\end{itemize}

\begin{definition}[Instruction-expression]
Инструкция-выражение — это выражение, завершённое \texttt{;}, например:
\begin{verbatim}
x = y + 1;
f();
\end{verbatim}
\end{definition}

\begin{remark}
Инициализация и присваивание — не одно и то же. Инициализация происходит при создании объекта, а присваивание — после его создания.
\end{remark}

\subsection{Логические выражения и control flow. \texttt{if}, \texttt{if} с инициализацией, \texttt{switch}, тернарный оператор}
\begin{definition}[Control flow]
\textit{Control flow} — это управление порядком выполнения инструкций в программе.
\end{definition}

Логические операции:
\begin{itemize}
    \item \texttt{\&\&} — logical and;
    \item \texttt{||} — logical or;
    \item \texttt{!} — logical not;
    \item сравнения: \texttt{==}, \texttt{!=}, \texttt{<}, \texttt{<=}, \texttt{>} , \texttt{>=}.
\end{itemize}

\texttt{if}:
\begin{verbatim}
if (cond) {
    ...
} else {
    ...
}
\end{verbatim}

\texttt{if} с инициализацией:
\begin{verbatim}
if (auto x = f(); x > 0) {
    ...
}
\end{verbatim}

\texttt{switch} удобен для выбора по целочисленному или перечислимому значению:
\begin{verbatim}
switch (n) {
case 1:
    ...
    break;
default:
    ...
}
\end{verbatim}

Тернарный оператор:
\begin{verbatim}
cond ? expr1 : expr2
\end{verbatim}

Он компактнее \texttt{if}, но удобен в основном для выбора значения, а не для сложной логики с несколькими инструкциями.

\subsection{Циклы: \texttt{while}, \texttt{for}, range-based for, \texttt{do..while}. \texttt{break}, \texttt{continue}, \texttt{return}}
\texttt{while} используется, когда число итераций заранее неизвестно:
\begin{verbatim}
while (cond) {
    ...
}
\end{verbatim}

\texttt{for} удобен для циклов со счётчиком:
\begin{verbatim}
for (int i = 0; i < n; ++i) {
    ...
}
\end{verbatim}

Range-based for применяют для обхода коллекции:
\begin{verbatim}
for (const auto& x : container) {
    ...
}
\end{verbatim}

\texttt{do..while} гарантирует хотя бы одну итерацию:
\begin{verbatim}
do {
    ...
} while (cond);
\end{verbatim}

\texttt{break} завершает цикл или \texttt{switch}, \texttt{continue} переходит к следующей итерации, \texttt{return} завершает функцию.

\begin{remark}
Range-based for можно переписать классическим \texttt{for} через итераторы, если коллекция поддерживает \texttt{begin()} и \texttt{end()}.
\end{remark}

\subsection{Краткие ответы на контрольные вопросы}
\begin{itemize}
    \item \textbf{Что такое инструкция?} Законченная синтаксическая конструкция языка, обычно отделяемая \texttt{;} или оформленная блоком.
    \item \textbf{Что такое описание?} Сообщение компилятору о сущности без обязательного создания.
    \item \textbf{Что такое определение?} Создание сущности, например объекта или функции.
    \item \textbf{Может ли быть только описание без определения?} Да, например \texttt{extern int x;}.
    \item \textbf{Может ли быть только определение без отдельного описания?} Да, потому что определение обычно само по себе является и описанием.
    \item \textbf{Можно ли использовать блок без заголовка?} Да, как отдельный локальный scope.
    \item \textbf{Как получить named block?} Через функцию.
    \item \textbf{Есть ли объект без типа?} В C++ нет: объект всегда имеет тип.
    \item \textbf{Какая разница между \texttt{short} и \texttt{int}?} Помимо размера, \texttt{int} обычно является базовым типом, к которому выполняются integer promotions; \texttt{short} часто приводится к \texttt{int} в выражениях.
    \item \textbf{Какая разница между \texttt{char} и \texttt{int}?} \texttt{char} — отдельный тип для символов и байтов; он может быть signed или unsigned и участвует в специальных правилах преобразования.
\end{itemize}
